% Unit 033 terjemahan Bahasa Indonesia dari chapter3.tex baris 163--345.
% Sumber: Methods of Algebra, Volume 2, commit
% 9a5803ff77dd3257484cb177f851a73770a59dd3 (CC BY 4.0).
% stable-unit-id: o014.aljabr2.chapter3.hom-complex-and-homotopy
% Cakupan: Bagian 3.2 lengkap; berhenti sebelum Bagian 3.3.

% segment-id: o014.li2.chapter3.hom-complex-and-homotopy.h001
\section{Kompleks \texorpdfstring{$\Hom$}{Hom} dan Homotopi}\label{sec:Hom-cplx}
\hypertarget{o014.aljabr2.chapter3.hom-complex-and-homotopy}{ }

% segment-id: o014.li2.chapter3.hom-complex-and-homotopy.p001
Sebelumnya telah didefinisikan $\Hom$ di antara dua kompleks. Bagian ini
menjelaskan cara meningkatkannya menjadi suatu kompleks, dan juga cara
meningkatkan komposisi morfisme menjadi sejenis perkalian pada kompleks.
Tetapkan kembali $\mathcal{A}$ sebagai kategori aditif. Jika $\Bbbk$ sebarang
gelanggang komutatif dan $\mathcal{A}$ kategori $\Bbbk$-linear, semua
pernyataan dalam bagian ini mempunyai versi $\Bbbk$-linear yang jelas.

% segment-id: o014.li2.chapter3.hom-complex-and-homotopy.p002
Ingat bahwa pengambilan himpunan $\Hom$ memberikan funktor
$\Hom:\mathcal{A}^{\opp}\times\mathcal{A}\to\cate{Ab}$. Notasi $\Hom$ tanpa
superskrip merujuk pada himpunan $\Hom$ dalam $\mathcal{A}$ atau
$\cate{C}(\mathcal{A})$; jika perlu, kategorinya dibedakan dengan subskrip.

% segment-id: o014.li2.chapter3.hom-complex-and-homotopy.d001
\begin{definisi}[Kompleks $\Hom$]\label{def:Hom-cplx}
	\index[sym1]{Hombullet@$\Hom^{\bullet}$}
	Misalkan $X,Y$ objek $\cate{C}(\mathcal{A})$. Untuk setiap $n\in\Z$,
	definisikan
% segment-id: o014.li2.chapter3.hom-complex-and-homotopy.q001
	\[ \Hom^n\left(X, Y \right) := \prod_{k \in \Z} \Hom_{\mathcal{A}}\left( X^k, Y^{k+n} \right); \]
	komposisi morfisme suku demi suku memberikan
% segment-id: o014.li2.chapter3.hom-complex-and-homotopy.q002
	\begin{equation}\label{eqn:Hom-cplx-multiplication}\begin{aligned}
		\Hom^n \left(Y, Z \right) \times \Hom^m\left( X, Y \right) & \longrightarrow \Hom^{n+m}\left(X , Z \right) \\
		(g, f) & \longmapsto gf := \left( g^{k+m} f^k \right)_{k \in \Z}.
	\end{aligned}\end{equation}
	Perhatikan bahwa $d_X=(d_X^k)_k\in\Hom^1(X,X)$, dan demikian pula untuk
	$d_Y$. Berdasarkan definisi ini, tetapkan
% segment-id: o014.li2.chapter3.hom-complex-and-homotopy.q003
	\begin{align*}
		d^n = d_{\Hom^\bullet(X, Y)}^n: \Hom^n\left(X, Y \right) & \longrightarrow \Hom^{n+1}\left(X, Y \right) \\
		f & \longmapsto d_Y f - (-1)^n f d_X.
	\end{align*}
\end{definisi}

% segment-id: o014.li2.chapter3.hom-complex-and-homotopy.p003
Operasi \eqref{eqn:Hom-cplx-multiplication} jelas asosiatif dan distributif
terhadap penjumlahan. Diberikan morfisme $u:\underline{X}\to X$ dan
$v:Y\to\underline{Y}$ dalam $\cate{C}(\mathcal{A})$, pandang keduanya
masing-masing sebagai unsur $\Hom^0(\underline{X},X)$ dan
$\Hom^0(Y,\underline{Y})$. Untuk setiap $n\in\Z$ diperoleh homomorfisme
% segment-id: o014.li2.chapter3.hom-complex-and-homotopy.q004
\begin{equation}\label{eqn:Hom-cplx-functoriality}\begin{aligned}
	\Hom^n(u, v): \Hom^n\left( X, Y\right) & \longrightarrow \Hom^n\left(\underline{X}, \underline{Y}\right) \\
	f & \longmapsto vfu.
\end{aligned}\end{equation}

% segment-id: o014.li2.chapter3.hom-complex-and-homotopy.dp001
\begin{definisiproposisi}
	\index{kompleks Hom@kompleks $\Hom$ ($\Hom$ complex)}
	Data $\left(\Hom^n(X,Y),d^n\right)_{n\in\Z}$ di atas memberikan objek
	dalam $\cate{C}(\cate{Ab})$, yang disebut \emph{kompleks $\Hom$}. Ketika
	$X$ dan $Y$ bervariasi, konstruksi ini bersama
	dengan \eqref{eqn:Hom-cplx-functoriality} memberikan funktor aditif
% segment-id: o014.li2.chapter3.hom-complex-and-homotopy.q005
	\[ \Hom^\bullet: \cate{C}(\mathcal{A})^{\opp} \times \cate{C}(\mathcal{A}) \to \cate{C}(\cate{Ab}). \]
\end{definisiproposisi}
% segment-id: o014.li2.chapter3.hom-complex-and-homotopy.pf001
\begin{bukti}
	Pertama-tama, verifikasi $d^{n+1}d^n=0$. Karena
	$d_Y^2=0=d_X^2$, perhitungan langsung memberikan
% segment-id: o014.li2.chapter3.hom-complex-and-homotopy.q006
	\begin{align*}
		d^{n+1} \left( d^n f \right) & = d_Y \left( d_Y f - (-1)^n f d_X \right) - (-1)^{n+1} \left( d_Y f - (-1)^n f d_X \right) d_X \\
		& = (-1)^{n+1} d_Y f d_X - (-1)^{n+1} d_Y f d_X = 0.
	\end{align*}

	Berikutnya, periksa funktorialitas. Masukkan $\Hom^\bullet(u,v)$ yang
	didefinisikan dalam \eqref{eqn:Hom-cplx-functoriality} ke dalam diagram%
	\footnote{Catatan penerjemah (O014-C033): sumber memberi label
	$\Hom^{n+1}(u,v)$ pada kedua anak panah horizontal. Karena anak panah atas
	menghubungkan suku berderajat $n$, target memulihkan labelnya menjadi
	$\Hom^n(u,v)$; anak panah bawah tetap berlabel $\Hom^{n+1}(u,v)$.}
% segment-id: o014.li2.chapter3.hom-complex-and-homotopy.g001
	\[\begin{tikzcd}
		\Hom^n\left( X^\bullet, Y^\bullet \right) \arrow[d, "{d^n}"'] \arrow[r, "{\Hom^n(u, v)}" inner sep=0.7em] & \Hom^n\left( \underline{X}^\bullet, \underline{Y}^\bullet \right) \arrow[d, "{\underline{d}^n}"] \\
		\Hom^{n+1}\left( X^\bullet, Y^\bullet \right) \arrow[r, "{\Hom^{n+1}(u, v)}"' inner sep=0.7em] & \Hom^{n+1}\left( \underline{X}^\bullet, \underline{Y}^\bullet \right) ;
	\end{tikzcd}\]
	Karena $d_Xu=ud_{\underline{X}}$ dan $vd_Y=d_{\underline{Y}}v$, diagram
	ini jelas komutatif.
\end{bukti}

% segment-id: o014.li2.chapter3.hom-complex-and-homotopy.p004
Selain itu, diferensial $d_{\Hom^\bullet(X,Y)}$ kompatibel dengan pergeseran.
Pembuktiannya jauh lebih mudah daripada pernyataannya.

% segment-id: o014.li2.chapter3.hom-complex-and-homotopy.le001
\begin{lema}\label{prop:Hom-cplx-d-translate}
	Misalkan $n,m\in\Z$ dan $f\in\Hom^n(X,Y)$. Identifikasikan $f$ dengan
	unsur $\underline{f}=(f^k)_k$ dari $\Hom^{n-m}(X,Y[m])$, atau dengan
	unsur $\overline{f}=(f^{k-m})_k$ dari $\Hom^{n-m}(X[-m],Y)$. Maka
% segment-id: o014.li2.chapter3.hom-complex-and-homotopy.q007
	\begin{align*}
		d^{n-m}_{\Hom^\bullet(X, Y[m])} \underline{f} & \xlongequal{\text{diidentifikasi dengan}} (-1)^m d^n_{\Hom^\bullet(X, Y)} f, \\ d^{n-m}_{\Hom^\bullet(X[-m], Y)} \overline{f} & \xlongequal{\text{diidentifikasi dengan}} d^n_{\Hom^\bullet(X, Y)} f.
	\end{align*}
	Akibatnya,
	$\Hom^\bullet(X,Y[m])=\Hom^\bullet(X,Y)[m]\simeq
	\Hom^\bullet(X[-m],Y)$.
\end{lema}
% segment-id: o014.li2.chapter3.hom-complex-and-homotopy.pf002
\begin{bukti}
	Sebagai contoh, untuk persamaan pertama,
% segment-id: o014.li2.chapter3.hom-complex-and-homotopy.q008
	\[ (d^{n-m}_{\Hom^\bullet(X, Y[m])} \underline{f})^k = d_{Y[m]}^{k+n-m} \underline{f}^k - (-1)^{n-m} \underline{f}^{k+1} d_X^k : X^k \to Y^{k+n+1}, \]
	dan ruas ini juga sama dengan
	$(-1)^m\left(d_Y^{k+n}f^k-(-1)^nf^{k+1}d_X^k\right)$. Argumen untuk
	$\overline{f}$ serupa: walaupun kasus $\overline{f}$ tidak mempunyai
	tanda $(-1)^m$, tetap ada isomorfisme kompleks $\Hom$; lihat catatan kaki
	pada Definisi~\ref{def:translation-functor}.
\end{bukti}

% segment-id: o014.li2.chapter3.hom-complex-and-homotopy.p005
Operasi perkalian yang didefinisikan dalam
\eqref{eqn:Hom-cplx-multiplication} juga memenuhi aturan Leibniz.

% segment-id: o014.li2.chapter3.hom-complex-and-homotopy.le002
\begin{lema}\label{prop:Hom-cplx-Leibniz}
	Misalkan $X,Y,Z$ kompleks dan ambil
	\[
		(g,f)\in\Hom^n(Y,Z)\times\Hom^m(X,Y).
	\]
	Untuk perkalian dalam \eqref{eqn:Hom-cplx-multiplication}, berlaku persamaan
% segment-id: o014.li2.chapter3.hom-complex-and-homotopy.q009
	\[
		\begin{aligned}
			d^{n+m}(gf)
				&= (d^n g)f \\
				&\quad + (-1)^n g(d^m f)
		\end{aligned}
	\]
	di $\Hom^{n+m+1}(X,Z)$.
\end{lema}
% segment-id: o014.li2.chapter3.hom-complex-and-homotopy.pf003
\begin{bukti}
	Lakukan perhitungan langsung; rinciannya diserahkan kepada pembaca.
\end{bukti}

% segment-id: o014.li2.chapter3.hom-complex-and-homotopy.le003
\begin{lema}\label{prop:Hom-as-cocycle}
	Misalkan $n\in\Z$ dan $f\in\Hom^n(X,Y)$. Maka $d^nf=0$ jika dan hanya
	jika $f\in\Hom(X,Y[n])$.
\end{lema}
% segment-id: o014.li2.chapter3.hom-complex-and-homotopy.pf004
\begin{bukti}
	Melalui Lema~\ref{prop:Hom-cplx-d-translate}, identifikasikan $f$ dengan
	unsur $\Hom^0(X,Y[n])$, sehingga masalahnya tereduksi pada kasus $n=0$.
	Jelas bahwa $d^0f=0\iff d_Yf=fd_X$.
\end{bukti}

% segment-id: o014.li2.chapter3.hom-complex-and-homotopy.d002
\begin{definisi}\label{def:homotopy}
	\index{homotopi (homotopy)}
	\index{homotopik nol (null-homotopic)}
	Misalkan $n\in\Z$. Pandang $\Hom(X,Y[n])$ sebagai subhimpunan
	$\Hom^n(X,Y)$.
% segment-id: o014.li2.chapter3.hom-complex-and-homotopy.l001
	\begin{compactitem}
% segment-id: o014.li2.chapter3.hom-complex-and-homotopy.i001
		\item Jika $f,g\in\Hom(X,Y[n])$, sedangkan
		$h\in\Hom^{n-1}(X,Y)$ (atau secara ekuivalen
		$h\in\Hom^{-1}(X,Y[n])$)%
		\footnote{Catatan penerjemah (O014-C034): sumber tidak menyatakan
		derajat $h$, padahal $d^{n-1}h$ menuntut
		$h\in\Hom^{n-1}(X,Y)$. Target mengikat tipenya secara eksplisit.}
		memenuhi $g-f=d^{n-1}h$, maka $h$ disebut
		suatu \emph{homotopi} dari $f$ ke $g$.
% segment-id: o014.li2.chapter3.hom-complex-and-homotopy.i002
		\item Untuk $f,g$ seperti di atas, jika terdapat suatu homotopi dari
		$f$ ke $g$, maka $f$ dan $g$ disebut homotopik. Morfisme $f$ yang
		homotopik dengan morfisme nol disebut \emph{homotopik nol}.
	\end{compactitem}

	Sebagai kasus khusus $n=0$, morfisme $f\in\Hom(X,Y)$ homotopik nol jika
	dan hanya jika terdapat keluarga morfisme $h^m:X^m\to Y^{m-1}$ sedemikian
	sehingga
% segment-id: o014.li2.chapter3.hom-complex-and-homotopy.q010
	\[ \forall m \in \Z, \quad f^m = d_Y^{m-1} h^m + h^{m+1} d_X^m . \]
\end{definisi}

% segment-id: o014.li2.chapter3.hom-complex-and-homotopy.p006
Relasi homotopi merupakan relasi ekuivalensi.
Lema~\ref{prop:Hom-as-cocycle}
memberikan persamaan penting berikut dalam $\cate{Ab}$ (atau
$\Bbbk\dcate{Mod}$):
% segment-id: o014.li2.chapter3.hom-complex-and-homotopy.q011
\[ \Hm^n\left(\Hom^\bullet(X, Y), d^\bullet \right) = \Hom\left(X, Y[n]\right) \big/ \{\text{morfisme homotopik nol} \}. \]

% segment-id: o014.li2.chapter3.hom-complex-and-homotopy.p007
Hasil berikut menunjukkan bahwa morfisme homotopik nol bukan hanya tertutup
terhadap penjumlahan, melainkan juga membentuk suatu ideal.

% segment-id: o014.li2.chapter3.hom-complex-and-homotopy.le004
\begin{lema}\label{prop:null-homotopic-composition}
	Misalkan $X\xrightarrow{f}Y[m]$ dan $Y\xrightarrow{g}Z[n]$ morfisme di
	antara kompleks. Jika $f$ atau $g$ homotopik nol, maka
	$g[m]\circ f:X\to Z[n+m]$ homotopik nol.
\end{lema}
% segment-id: o014.li2.chapter3.hom-complex-and-homotopy.pf005
\begin{bukti}
	Superskrip pada $d$ dihilangkan di bawah ini. Misalkan $f=dh$, dengan
	$h\in\Hom^{-1}(X,Y[m])$. Dengan memahami komposisi di dalam
	$\Hom^\bullet$, kita dapat menyingkat $g[m]\circ f$ menjadi $gf$.
	Lema~\ref{prop:Hom-cplx-Leibniz} mengimplikasikan
	$d(gh)=(dg)h+(-1)^ng(dh)=(-1)^ngf$, sehingga $gf$ homotopik nol. Kasus
	$g=dh$ serupa.
\end{bukti}

% segment-id: o014.li2.chapter3.hom-complex-and-homotopy.p008
Dengan menghilangkan subskrip pada $d_{\Hom^\bullet}$,
Lema~\ref{prop:Hom-cplx-Leibniz} dan
Lema~\ref{prop:null-homotopic-composition} menunjukkan bahwa perkalian pada
$\Hom^\bullet$ dalam \eqref{eqn:Hom-cplx-multiplication} menginduksi
% segment-id: o014.li2.chapter3.hom-complex-and-homotopy.q012
\begin{equation}\label{eqn:homotopy-cat-cmposition}
	\Hm^n\left(\Hom^\bullet(Y, Z) \right) \times \Hm^m\left(\Hom^\bullet(X, Y) \right) \to \Hm^{n+m}\left( \Hom^\bullet(X, Z) \right) ,
\end{equation}
dan operasi biner \eqref{eqn:homotopy-cat-cmposition} mewarisi sifat asosiatif
dan distributif terhadap penjumlahan dari
\eqref{eqn:Hom-cplx-multiplication}.

% segment-id: o014.li2.chapter3.hom-complex-and-homotopy.d003
\begin{definisi}\label{def:cplx-homotopy-cat}
	\index[sym1]{KA@$\cate{K}(\mathcal{A})$}
	Definisikan kategori $\cate{K}(\mathcal{A})$ dari
	$\cate{C}(\mathcal{A})$ dengan menetapkan
% segment-id: o014.li2.chapter3.hom-complex-and-homotopy.q013
	\begin{align*}
		\Obj\left( \cate{K}(\mathcal{A}) \right) & := \Obj\left( \cate{C}(\mathcal{A}) \right), \\
		\Hom_{\cate{K}(\mathcal{A})}\left( X, Y \right) & := \Hm^0\left( \Hom^\bullet\left( X, Y \right), d^\bullet_{\Hom} \right),
	\end{align*}
	dengan komposisi morfisme ditentukan oleh
	\eqref{eqn:homotopy-cat-cmposition} untuk $n=m=0$.
\end{definisi}

% segment-id: o014.li2.chapter3.hom-complex-and-homotopy.p009
Perhatikan bahwa $\cate{K}(\mathcal{A})$ merupakan kategori aditif:
penjumlahan pada himpunan morfisme sudah jelas, objek nol berasal dari objek
nol dalam $\cate{C}(\mathcal{A})$, sedangkan keberadaan produk (atau
koproduk) berasal dari sifat yang lebih umum berikut. Terdapat isomorfisme
alami
% segment-id: o014.li2.chapter3.hom-complex-and-homotopy.q014
\begin{align*}
	\Hom_{\cate{K}(\mathcal{A})}\left( X, \prod_{i \in I} Y_i \right) & \simeq \prod_{i \in I} \Hom_{\cate{K}(\mathcal{A})}(X, Y_i) \\
	\text{atau}\quad \Hom_{\cate{K}(\mathcal{A})}\left( \coprod_{i \in I} X_i, Y \right) & \simeq \prod_{i \in I} \Hom_{\cate{K}(\mathcal{A})}(X_i, Y),
\end{align*}
dengan $I$ sebarang himpunan, asalkan $\prod_{i\in I}Y_i$ (atau
$\coprod_iX_i$) ada pada tingkat $\cate{C}(\mathcal{A})$. Karena pengambilan
produk dan pengambilan $\Hm^0$ berkomutasi dalam $\cate{C}(\cate{Ab})$,
semuanya tereduksi pada isomorfisme kompleks $\Hom$ yang jelas:
% segment-id: o014.li2.chapter3.hom-complex-and-homotopy.q015
\begin{align*}
	\Hom^\bullet\left(X, \prod_i Y_i\right) & \simeq \prod_i \Hom^\bullet(X, Y_i) \\
	\text{atau}\quad \Hom^\bullet\left(\coprod_i X_i, Y\right) & \simeq \prod_i \Hom^\bullet(X_i, Y).
\end{align*}

% segment-id: o014.li2.chapter3.hom-complex-and-homotopy.p010
Funktor pergeseran $[n]$ menginduksi automorfisme kategori aditif
$\cate{K}(\mathcal{A})$, yang tetap dinotasikan dengan $[n]$. Funktor aditif
$\cate{C}(\mathcal{A})\to\cate{K}(\mathcal{A})$ merupakan identitas pada
himpunan objek dan pemetaan kuosien pada himpunan morfisme. Sifat universal
berikut sudah jelas.

% segment-id: o014.li2.chapter3.hom-complex-and-homotopy.pr001
\begin{proposisi}\label{prop:homotopy-category-universal-property}
	Jika $\mathcal{B}$ kategori-$\cate{Ab}$ dan funktor aditif
	$F:\cate{C}(\mathcal{A})\to\mathcal{B}$ memetakan setiap morfisme
	homotopik nol ke morfisme nol, maka $F$ secara unik berfaktor melalui
	$\cate{C}(\mathcal{A})\to\cate{K}(\mathcal{A})$.
\end{proposisi}

% segment-id: o014.li2.chapter3.hom-complex-and-homotopy.p011
Latihan pada bab berikutnya akan menunjukkan bahwa $\cate{K}(\mathcal{A})$
pada umumnya bukan kategori abelian.

% segment-id: o014.li2.chapter3.hom-complex-and-homotopy.r001
\begin{catatan}\index{kompleks Hom}\index[sym1]{Homlbullet@$\Hom_\bullet$}
	\label{rem:Hom-chain-cplx}
	Untuk kompleks rantai (Catatan~\ref{rem:cochain-vs-chain}) juga terdapat
	$\Hom_\bullet$ yang bersesuaian:
% segment-id: o014.li2.chapter3.hom-complex-and-homotopy.q016
	\begin{gather*}
		\Hom_n(X, Y) := \prod_{k \in \Z} \Hom\left(X_k, Y_{k+n} \right), \\
		d_n f = d_Y f - (-1)^n f d_X, \quad f \in \Hom_n(X, Y).
	\end{gather*}
	Sesungguhnya, gunakan cara dalam Catatan~\ref{rem:cochain-vs-chain} dengan
	menetapkan $X^n=X_{-n}$, $d^n=d_{-n}$, dan seterusnya. Maka
	$\left(\Hom_n(X,Y),d_n\right)_n$ tepat merupakan kompleks rantai yang
	bersesuaian dengan $\left(\Hom^n(X,Y),d^n\right)_n$:
% segment-id: o014.li2.chapter3.hom-complex-and-homotopy.q017
	\begin{align*}
		\Hom_n\left(X, Y \right) & = \prod_{k \in \Z} \Hom\left( X_k, Y_{n+k} \right) \xlongequal{h = -k} \prod_{h \in \Z} \Hom\left(X^h, Y^{h-n} \right) \\
		& = \Hom^{-n} \left(X, Y\right),
	\end{align*}
	dan hubungan antara $d_n$ dan $d^{-n}$ juga serupa. Tentu saja, untuk
	kompleks rantai juga terdapat versi $\cate{K}(\mathcal{A})$ yang
	bersesuaian.
\end{catatan}

% segment-id: o014.li2.chapter3.hom-complex-and-homotopy.p012
Diberikan funktor aditif $F:\mathcal{A}\to\mathcal{B}$ di antara kategori
aditif, untuk setiap $X,Y\in\Obj(\cate{C}(\mathcal{A}))$, penerapan $F$ suku
demi suku memberikan morfisme dalam $\cate{C}(\cate{Ab})$
% segment-id: o014.li2.chapter3.hom-complex-and-homotopy.q018
\begin{align*}
	\Hom^\bullet\left(X, Y\right) & \to \Hom^\bullet\left(\cate{C}F(X), \cate{C}F(Y) \right) \\
	(f^k)_{k \in \Z}\in\Hom^r(X,Y) & \mapsto
	(F(f^k))_{k \in \Z}\in
	\Hom^r\left(\cate{C}F(X),\cate{C}F(Y)\right) \qquad (r\in\Z).
\end{align*}
Dengan mengambil $\Hm^0$ pada kedua ruas, diperoleh
\[
	\Hom_{\cate{K}(\mathcal{A})}(X,Y)\longrightarrow
	\Hom_{\cate{K}(\mathcal{B})}
	\bigl(\cate{C}F(X),\cate{C}F(Y)\bigr).
\]
Dengan demikian,
diperoleh funktor $\cate{K}F:\cate{K}(\mathcal{A})\to
\cate{K}(\mathcal{B})$ yang membuat diagram berikut komutatif:
% segment-id: o014.li2.chapter3.hom-complex-and-homotopy.g002
\begin{equation}\label{eqn:KF}\begin{tikzcd}
	\cate{C}(\mathcal{A}) \arrow[r, "{\cate{C}F}"] \arrow[d] & \cate{C}(\mathcal{B}) \arrow[d] \\
	\cate{K}(\mathcal{A}) \arrow[r, "{\cate{K}F}"'] & \cate{K}(\mathcal{B}).
\end{tikzcd}\end{equation}
Baris atasnya adalah funktor pada Proposisi~\ref{prop:CA-functoriality}.
Konstruksi-konstruksi ini kompatibel dengan pasangan adjoin.
\index[sym1]{KF@$\cate{K}F$}

% segment-id: o014.li2.chapter3.hom-complex-and-homotopy.pr002
\begin{proposisi}\label{prop:KF-adjoint}
	Tinjau pasangan adjoin $(F,G,\varphi)$, dengan
% segment-id: o014.li2.chapter3.hom-complex-and-homotopy.g003
	$\begin{tikzcd}
		F: \mathcal{A} \arrow[shift left, r] & \mathcal{B} \arrow[shift left, l] : G
	\end{tikzcd}$
	funktor-funktor aditif di antara kategori aditif. Maka
% segment-id: o014.li2.chapter3.hom-complex-and-homotopy.g004
	\[\begin{tikzcd}
		\cate{K}F: \cate{K}(\mathcal{A}) \arrow[shift left, r] & \cate{K}(\mathcal{B}) \arrow[shift left, l] : \cate{K}G
	\end{tikzcd}\]
	juga secara alami merupakan pasangan adjoin.
\end{proposisi}
% segment-id: o014.li2.chapter3.hom-complex-and-homotopy.pf006
\begin{bukti}
	Data pasangan adjoin $\varphi$ merupakan keluarga bijeksi
	$\varphi_{A,B}:\Hom_{\mathcal{B}}(FA,B)\simeq
	\Hom_{\mathcal{A}}(A,GB)$ yang funktorial dalam $A$ dan $B$.
	Proposisi~\ref{prop:automatic-k-morphism} memastikan bahwa bijeksi-bijeksi
	ini linear, sehingga menginduksi isomorfisme kanonik kompleks
% segment-id: o014.li2.chapter3.hom-complex-and-homotopy.q019
	\[ \Hom^\bullet\left( \cate{C}F(X), Y \right) \simeq \Hom^\bullet\left( X, \cate{C}G(Y) \right). \]
	Pada derajat $r$, isomorfisme ini memetakan
	\[
		(f^k)_{k\in\Z}\in\Hom^r(\cate{C}F(X),Y)
		\longmapsto
		(\varphi_{X^k,Y^{k+r}}(f^k))_{k\in\Z}
		\in\Hom^r(X,\cate{C}G(Y)).
	\]%
	\footnote{Catatan penerjemah (O014-C035): sumber menuliskan komponen
	sebagai $\varphi_{A,B}(f^n)$ tanpa mengikat $A$ dan $B$ pada objek
	berderajat yang menjadi sumber dan target $f^n$. Target memakai indeks
	derajat kompleks $r$ dan indeks keluarga $k$, sehingga komponennya
	tertipe sebagai $\varphi_{X^k,Y^{k+r}}(f^k)$.}
	Berdasarkan \eqref{eqn:KF}, pengambilan $\Hm^0$
	pada kedua ruas menjadikan $\cate{K}F$ adjoin kiri dari $\cate{K}G$.
\end{bukti}
